\subsection{El término local de la fórmula canónica}\label{sec:P04}
Se toma $n$ suficientemente grande, en particular $n\ge8M$.
Se distingue una consecuencia de la Sección~\ref{sec:P02} de otra que requiere
la comparación con el perfil continuo (Secciones~\ref{sec:P05} y~\ref{sec:P06}).
El término local es
\[
 L_n=-\frac12\sum_{ij}
       [\log(1-4\pi_{ij})-4\log(1-\pi_{ij})].       \tag{LC1}
\]
Para $0\le t\le1/8$, las series escalares convergentes dan
\[
 -\tfrac12[\log(1-4t)-4\log(1-t)]
   =\frac12\sum_{m\ge2}\frac{4^m-4}{m}t^m.
\]
El término cuadrático es $3t^2$ y los restantes son positivos.
Su suma está acotada por
\[
 \frac12\sum_{m\ge3}(4t)^m
       =\frac{32t^3}{1-4t}\le64t^3.
\]
Con $\pi\le M/n$ y $\sum\pi=n$ resulta
\[
 0\le L_n-3\sum_{ij}\pi_{ij}^2
       \le64\sum\pi_{ij}^3\le\frac{64M^2}{n}.      \tag{LC2}
\]
La estimación (LC2) es una expansión local de escala $\pi=O(1/n)$.

\paragraph{Paso del término local al perfil continuo.}
Si $\nu_1$ es el perfil continuo acotado y Lipschitz y
\[
 \max_{ij}|n\pi_{ij}-\nu_1((i-1/2)/n,(j-1/2)/n)|=O(1/n),
                                                        \tag{LC3}
\]
entonces
\[
 \sum_{ij}\pi_{ij}^2
       =\int_{[0,1]^2}\nu_1(r,c)^2\,dr\,dc+O(1/n).
\]
En efecto, el error entre los cuadrados de los perfiles
es $O(1/n)$ uniformemente, y la cuadratura bidimensional
de una función Lipschitz tiene error $O(1/n)$.
Por (LC2),
\[
 L_n=3\int_{[0,1]^2}\nu_1^2+O(1/n).               \tag{LC4}
\]
Éste es sólo uno de los sumandos de $\log Q_{\rm pred}$; las contribuciones de la entropía y
del pseudodeterminante se estudian en la Sección~\ref{sec:constante}.

Las Secciones~\ref{sec:P05} y~\ref{sec:P06} prueban (LC3) para el maximizador canónico;
por tanto (LC4) se aplica a él.
